\mxunit[13]{terminale-c13}{Représentations paramétriques et équations cartésiennes}\label{ch:12}

\mxsubtitle{S65, S67, S70}

\begin{keypoint}[Objectifs]Décrire une droite et un plan ; déterminer des intersections ; choisir une méthode pour une projection.\end{keypoint}

\section{Notes de cours}
Dans tout ce chapitre, les coordonnées sont données dans un repère orthonormé de l'espace. Cette hypothèse est nécessaire pour identifier un vecteur normal à partir des coefficients d'une équation et pour calculer les distances.
\subsection{Droites paramétrées}
La droite passant par $A\pt{x_A;y_A;z_A}$ de vecteur directeur non nul $\pt{a;b;c}$ admet la représentation
\[d:\quad\begin{cases}x=x_A+at\\y=y_A+bt\\z=z_A+ct\end{cases}\quad(t\in\R).\]
Un point appartient à la droite s'il existe \emph{un même} réel $t$ satisfaisant les trois égalités. Une droite admet de nombreuses représentations : changement du point de départ ou multiplication du directeur par un réel non nul.
\paragraph{Exemple.} $\pt{x;y;z}=\pt{1;2;0}+t\pt{2;-1;3}$ décrit la même droite que $\pt{3;1;3}+s\pt{4;-2;6}$ ; le lien est $t=1+2s$.
\subsection{Plans cartésiens}
\begin{keypoint}[Équation d'un plan]
Si $\pt{a;b;c}\ne\pt{0;0;0}$, l'ensemble $ax+by+cz+d=0$ est un plan de vecteur normal $\pt{a;b;c}$. Le plan de normal $\vec n$ passant par $A$ se construit en écrivant $\vect{AM}\cdot\vec n=0$.
\end{keypoint}
\paragraph{Démonstration.} La condition $\vect{AM}\cdot\vec n=0$ s'écrit $a(x-x_A)+b(y-y_A)+c(z-z_A)=0$. En développant, $d=-ax_A-by_A-cz_A$. Réciproquement, tout point satisfaisant cette équation vérifie la condition d'orthogonalité, donc appartient au plan.
\paragraph{Exemple.} Le plan de normal $\pt{2;-1;3}$ passant par $A\pt{1;2;0}$ est $2(x-1)-(y-2)+3z=0$, soit $2x-y+3z=0$.
\subsection{Intersections}
Pour une droite et un plan, remplacer $x,y,z$ dans l'équation du plan par leurs expressions paramétriques. Une solution en $t$ donne une intersection ponctuelle ; une identité donne une droite incluse ; une contradiction donne une droite parallèle disjointe.
Pour deux plans, des normales non colinéaires donnent une intersection suivant une droite ; résoudre le système à deux équations en choisissant une variable comme paramètre. Si les normales sont colinéaires, les plans sont parallèles : distinguer confondus et disjoints.
Pour deux droites, utiliser des paramètres $s,t$ indépendants et résoudre les trois équations. Sans solution et avec directeurs non colinéaires, les droites sont non coplanaires.
\subsection{Projection sur un plan}
Si $\mathcal P:ax+by+cz+d=0$, chercher le projeté de $M$ sur la droite normale passant par $M$ : $H=M+t\pt{a;b;c}$. Remplacer dans l'équation du plan pour déterminer $t$, puis calculer $MH$.
On en déduit, sans en faire une formule à apprendre isolément,
\[d(M,\mathcal P)=\frac{|ax_M+by_M+cz_M+d|}{\sqrt{a^2+b^2+c^2}}.\]
\paragraph{Vérification indispensable.} Après le calcul d'un projeté, contrôler simultanément $H\in\mathcal P$ et la colinéarité de $\vect{MH}$ avec le normal.


\mxnewpage
\section{Activité de recherche}\textbf{Trois droites face à un plan --- 30 min.}\par Dans un repère orthonormé de l'espace, on considère le plan $\mathcal P$ d'équation $x+y+z=3$. Dans chaque question, le paramètre $t$ parcourt $\R$.
\begin{enumerate}\item La droite $d_1$ admet la représentation $\pt{x;y;z}=\pt{t;t;t}$. Déterminer $d_1\cap\mathcal P$.
\item La droite $d_2$ admet la représentation $\pt{x;y;z}=\pt{3+t;-t;0}$. Déterminer $d_2\cap\mathcal P$.
\item La droite $d_3$ admet la représentation $\pt{x;y;z}=\pt{t;-t;0}$. Déterminer $d_3\cap\mathcal P$.
\item Pour chaque droite, préciser si l'équation obtenue après substitution dans celle du plan a une solution unique, aucune solution ou tout réel pour solution. Relier ces cas aux positions géométriques trouvées.\end{enumerate}

\mxnewpage
\section{Exercices progressifs}

\mxexerciseitem{1}{Entraînement}Dans un repère de l'espace, on considère $A\pt{1;-1;2}$ et $B\pt{3;0;1}$. Donner une représentation paramétrique de la droite $(AB)$.

\mxexerciseitem{2}{Entraînement}Dans un repère orthonormé, déterminer une équation cartésienne du plan passant par $A\pt{1;2;3}$ et de vecteur normal $\vec n=\pt{1;0;-2}$.

\mxexerciseitem{3}{Entraînement}Dans un repère de l'espace, la droite $d$ est donnée par $\pt{x;y;z}=\pt{t;1+t;2-t}$, $t\in\R$, et le plan $\mathcal P$ par $x+2y+z=8$. Déterminer leur intersection.

\mxexerciseitem{4}{Entraînement}Dans un repère de l'espace, déterminer l'intersection des plans d'équations $x+y+z=1$ et $x-y=0$. Donner une représentation paramétrique si l'intersection est une droite.

\mxexerciseitem{5}{Entraînement}Dans un repère de l'espace, étudier la position relative des plans d'équations $2x-y+z=3$ et $4x-2y+2z=5$. Justifier.

\section{Travail dans le manuel ou une annale}\textbf{Baccalauréat, Métropole, 17 juin 2025 : exercice 3, affirmations 1 à 4 (p. 5--6).} À la maison après S70. Justifier chaque réponse ; pour une affirmation fausse, expliquer précisément ce qui la réfute. Le sujet est accessible ici : \href{https://www.education.gouv.fr/sites/default/files/2025-06/baccalaur-at-g-n-ral-2025-math-matiques-440856.pdf}{annale officielle}. Utiliser les données et les consignes de cette annale.

\mxnewpage
\section{Devoir à la maison}\textbf{DM12 --- Projection et intersection --- 55 min.}\par Dans un repère orthonormé de l'espace, on considère le point $M\pt{3;0;0}$ et les plans $\mathcal P:x+2y+2z=9$ et $\mathcal Q:x-y=0$.
\begin{enumerate}\item Donner les coordonnées d'un vecteur normal non nul à $\mathcal P$.
\item Déterminer les coordonnées du projeté orthogonal $H$ de $M$ sur $\mathcal P$, puis calculer la distance de $M$ à $\mathcal P$.
\item Montrer que $\mathcal P\cap\mathcal Q$ est une droite et en donner une représentation paramétrique.
\item Le point $H$ appartient-il à cette droite ? Justifier.
\item Déterminer les coordonnées du point $K$ de cette droite dont la cote est nulle. Calculer $MK$, puis comparer $MK$ et $MH$.\end{enumerate}

\ifmxteacher\mxnewpage\section{Conduite et corrigés}\begin{teachernote}[Repères pour le professeur]Utiliser deux paramètres différents pour comparer deux droites. Vérifier les trois coordonnées, même si deux suffisent à calculer des candidats. Une équation cartésienne seule décrit en général un plan dans l'espace ; elle ne décrit pas une droite. Réinvestir le projeté du chapitre~\ref{ch:11} pour donner du sens à la résolution de systèmes.\end{teachernote}

\subsection{Corrigé de l'activité}\begin{enumerate}
\item Un point de $d_1$ a les coordonnées $\pt{t;t;t}$. Il appartient au plan si et seulement si $3t=3$, soit $t=1$. L'intersection est donc l'unique point $\pt{1;1;1}$.
\item Pour un point de $d_2$, $x+y+z=3+t-t+0=3$ quel que soit $t$. Tous les points de la droite sont dans le plan : $d_2\subset\mathcal P$.
\item Pour un point de $d_3$, $x+y+z=t-t+0=0$. L'équation $0=3$ n'a aucune solution : la droite est disjointe du plan. Son directeur $\pt{1;-1;0}$ est orthogonal au normal $\pt{1;1;1}$ ; elle est parallèle au plan.
\item Après substitution, on obtient respectivement une équation avec une seule solution, une identité vraie pour tout paramètre, et une contradiction. Ces trois cas correspondent à une intersection ponctuelle, une inclusion et une absence d'intersection.
\end{enumerate}

\paragraph{E1.}$\vect{AB}=\pt{2;1;-1}\ne\vec0$. La droite a donc pour représentation $\pt{x;y;z}=\pt{1;-1;2}+t\pt{2;1;-1}$, $t\in\R$. Les valeurs $t=0$ et $t=1$ redonnent $A$ et $B$.

\paragraph{E2.}Pour $A\pt{1;2;3}$, la condition $\vect{AM}\cdot\pt{1;0;-2}=0$ est $(x-1)-2(z-3)=0$. Après développement, le plan a pour équation $x-2z+5=0$. Le point $A$ la vérifie.

\paragraph{E3.}Après substitution, $t+2(1+t)+(2-t)=8$, soit $2t+4=8$, donc $t=2$. L'unique intersection est $\pt{2;3;0}$ ; on vérifie $2+2\times3+0=8$.

\paragraph{E4.}La seconde équation donne $x=y$. Posons $x=y=t$ ; la première devient $2t+z=1$. L'ensemble des solutions est $\pt{x;y;z}=\pt{t;t;1-2t}=\pt{0;0;1}+t\pt{1;1;-2}$, une droite.

\paragraph{E5.}Les normaux sont colinéaires, le second étant le double du premier. Mais doubler $2x-y+z=3$ donne $4x-2y+2z=6$, incompatible avec le second membre $5$. Les plans sont donc strictement parallèles, et non sécants.

\mxneed{30}

\subsection{Corrigé du DM12}\begin{enumerate}
\item Les coefficients de $x,y,z$ donnent le normal non nul $\vec n=\pt{1;2;2}$ dans le repère orthonormé.
\item Cherchons $H$ sur la normale passant par $M$ : $H=\pt{3+t;2t;2t}$. L'appartenance à $\mathcal P$ donne $3+t+4t+4t=9$, donc $t=\frac{2}{3}$. Ainsi $H=\pt{\frac{11}{3};\frac{4}{3};\frac{4}{3}}$. On vérifie $\frac{11}{3}+\frac{8}{3}+\frac{8}{3}=9$ et $\vect{MH}=(\frac{2}{3})\vec{n}$. Donc $H$ est bien le projeté, et $MH=(\frac{2}{3})\sqrt{1+4+4}=2$.
\item Les normaux $\pt{1;2;2}$ et $\pt{1;-1;0}$ ne sont pas colinéaires : les plans se coupent suivant une droite. Comme $x=y$, posons $x=y=s$. La première équation donne $3s+2z=9$, soit $z=\frac{9-3s}{2}$. L'intersection est la droite $\pt{x;y;z}=\pt{s;s;\frac{9-3s}{2}}$, $s\in\R$.
\item Le point $H$ a $x_H=\frac{11}{3}\ne\frac{4}{3}=y_H$, donc il n'appartient pas à $\mathcal Q$ ni à la droite d'intersection.
\item La cote est la coordonnée $z$. Sa nullité impose $9-3s=0$, donc $s=3$ et $K=\pt{3;3;0}$. Alors $MK=3>2=MH$. Ce résultat est cohérent avec la propriété du projeté : parmi tous les points du plan, $H$ est le plus proche de $M$.
\end{enumerate}

\subsection{Corrigé du travail dans le manuel ou une annale}\begin{enumerate}
\item Vraie : la droite passe par $B$ et son vecteur directeur $\pt{-2;-1;3}$ vaut $-\frac12\vect{AB}$.
\item Fausse : $\vec n\cdot\vect{OB}=15-4-1=10\ne0$.
\item Fausse : pour $k=-4$ et $s=\frac52$, les deux représentations donnent $\pt{11;12;-14}$. Les droites se coupent, donc sont coplanaires.
\item Vraie : $d(C,\mathcal P)=\frac{|2-(-1)+2+1|}{\sqrt{1+1+1}}=2\sqrt3$.
\end{enumerate}

\fi

